\documentclass[11pt,a4paper]{article}

\usepackage[margin=2.4cm]{geometry}
\usepackage[T1]{fontenc}
\usepackage[utf8]{inputenc}
\usepackage[french]{babel}
\usepackage{lmodern}
\usepackage{enumitem}
\usepackage{xcolor}
\usepackage[colorlinks=true,urlcolor=blue,linkcolor=black,citecolor=black]{hyperref}
\usepackage{parskip}

\setlist{leftmargin=1.5em,itemsep=2pt,topsep=4pt}

\title{\textbf{CityVision --- Plan prévisionnel}\\[2pt]
\large Preuve de concept : segmentation sémantique de scènes urbaines\\
avec un modèle récent (EoMT)}
\author{Darya Filatova}
\date{}

\begin{document}
\sloppy
\maketitle

% =====================================================================
\section*{Dataset retenu}
% =====================================================================

\textbf{Cityscapes} \cite{cordts2016} : photos de scènes urbaines prises depuis
un véhicule dans des villes allemandes et voisines, en $2048\times1024$ pixels,
annotées pixel par pixel.

\begin{itemize}
  \item \textbf{Découpage :} 2\,975 images d'entraînement (18 villes), 500 de
        validation (3 villes) et 1\,525 de test, dont les labels ne sont pas
        publics.
  \item \textbf{Classes :} pour le projet, les classes Cityscapes sont
        regroupées en \textbf{9 classes} (fond, route, bâtiment, végétation,
        ciel, personne, voiture, panneau, vélo).
  \item \textbf{Résolution :} images ramenées à $512\times1024$.
  \item \textbf{Déséquilibre :} le jeu est très déséquilibré. Route, fond et
        bâtiments représentent environ 74\,\% des pixels ; personnes, panneaux
        et vélos moins de 2\,\% au total, alors que les panneaux apparaissent
        dans 94\,\% des images.
\end{itemize}

C'est le jeu de données de référence du projet : le modèle actuel (baseline) y
a été entraîné.

% =====================================================================
\section*{Modèle envisagé}
% =====================================================================

\textbf{Algorithme : EoMT} (\emph{Encoder-only Mask Transformer},
Kerssies et al., CVPR~2025) \cite{kerssies2025vitsecretlyimagesegmentation},
avec \textbf{Mask2Former} (Cheng et al., CVPR~2022) \cite{cheng2022} comme étape
intermédiaire.

\subsection*{Objectif et contexte}

La segmentation sémantique, c'est-à-dire attribuer une classe à chaque pixel
d'une image. C'est la base de la perception visuelle des véhicules autonomes et
des systèmes d'aide à la conduite, où il faut savoir précisément où sont la
route, les piétons ou les panneaux, idéalement en temps réel.

\subsection*{Arguments justifiant le choix}

\begin{itemize}
  \item \textbf{Un changement de paradigme.} La baseline (U-Net + ResNet50)
        classe chaque pixel indépendamment. Mask2Former puis EoMT prédisent un
        \textbf{ensemble de masques, chacun associé à une classe}. Cette
        formulation a donné l'état de l'art à sa publication sur COCO, ADE20K
        et Cityscapes, en segmentation sémantique, d'instance et panoptique
        \cite{cheng2022}.
  \item \textbf{Un gain attendu sur les petits objets.} Chaque masque couvre un
        objet entier, ce qui devrait aider sur les classes rares de notre jeu
        (personne, panneau, vélo), là où la baseline est la plus faible.
  \item \textbf{La simplicité et la vitesse d'EoMT.} EoMT montre qu'un Vision
        Transformer simple, pré-entraîné à grande échelle (DINOv2), suffit pour
        segmenter, sans les décodeurs spécialisés de Mask2Former. Les auteurs
        annoncent une précision comparable aux méthodes de l'état de l'art,
        jusqu'à \textbf{4$\times$ plus rapide} avec un ViT-L
        \cite{kerssies2025vitsecretlyimagesegmentation}. C'est un critère
        important pour une utilisation embarquée ou via une API.
  \item \textbf{La faisabilité.} Le code officiel est publié sous licence MIT,
        et les deux modèles sont intégrés à Hugging Face Transformers avec des
        poids pré-entraînés.
\end{itemize}

% =====================================================================
\section*{Références bibliographiques}
% =====================================================================

\renewcommand{\refname}{\vspace{-2.2em}}
\begin{thebibliography}{9}
\bibitem{cheng2022} B.~Cheng, I.~Misra, A.~G.~Schwing, A.~Kirillov,
  R.~Girdhar, \emph{Masked-attention Mask Transformer for Universal Image
  Segmentation}, CVPR 2022. arXiv:2112.01527,
  \url{https://arxiv.org/abs/2112.01527}.\\
  Cette référence introduit la classification de masques avec attention
  masquée, base de la méthode.
\bibitem{kerssies2025vitsecretlyimagesegmentation} T.~Kerssies, N.~Cavagnero,
  A.~Hermans, N.~Norouzi, G.~Averta, B.~Leibe, G.~Dubbelman, D.~de~Geus,
  \emph{Your ViT is Secretly an Image Segmentation Model}, 2025.
  arXiv:2503.19108 [cs.CV], \url{https://arxiv.org/abs/2503.19108}.
  Code : \url{https://github.com/tue-mps/eomt} (licence MIT).\\
  C'est l'algorithme mis en œuvre.
\bibitem{cordts2016} M.~Cordts et al., \emph{The Cityscapes Dataset for
  Semantic Urban Scene Understanding}, CVPR 2016. arXiv:1604.01685,
  \url{https://arxiv.org/abs/1604.01685}.\\
  C'est la référence du jeu de données et de son protocole d'évaluation.
\end{thebibliography}

\subsection*{Réutilisation de code}

\begin{itemize}
  \item implémentations Hugging Face Transformers
        (\texttt{Mask2Former\allowbreak For\allowbreak Universal\allowbreak Segmentation},
        \texttt{Eomt\allowbreak For\allowbreak Universal\allowbreak Segmentation}) ;
  \item dépôt officiel EoMT, dont je reprends la recette d'entraînement :
        décroissance du taux d'apprentissage par couche, \emph{warm-up} en deux
        temps, \emph{mask annealing}.
\end{itemize}

L'application est originale : le jeu de données (Cityscapes ramené à
9 classes, en $512\times1024$), le pipeline de données, l'évaluation et
l'interface sont propres au projet.

% =====================================================================
\section*{Démarche de test du nouvel algorithme (preuve de concept)}
% =====================================================================

\subsection*{Baseline}

\begin{itemize}
  \item le modèle actuel du projet, \textbf{ResNet50-UNet} (encodeur
        pré-entraîné sur ImageNet) ;
  \item entraîné 20 époques avec une perte cross-entropie + Dice pondérée ;
  \item mesuré par la \textbf{mIoU} (moyenne sur les classes de l'IoU,
        intersection sur union) sur les 500 images de validation.
\end{itemize}

\subsection*{Mise en œuvre}

\begin{enumerate}
  \item \textbf{Mask2Former (Swin-Tiny, pré-entraîné sur ADE20K), puis EoMT
        (ViT-B DINOv2, pré-entraîné sur COCO)}, ré-entraînés sur nos 9 classes
        \textbf{dans les mêmes conditions que la baseline} : mêmes images, même
        résolution, même métrique, même budget (20 époques, batch effectif
        de 4, arrêt anticipé). Chaque modèle garde sa recette d'entraînement
        publiée. Aucun point de départ n'a vu Cityscapes.
  \item \textbf{Comparaison :}
        \begin{itemize}
          \item mIoU et IoU par classe, en particulier sur les petites classes ;
          \item temps d'inférence par image et nombre de paramètres ;
          \item un \textbf{bootstrap apparié} sur les images de validation,
                pour vérifier que l'écart n'est pas dû au hasard (intervalle de
                confiance à 95\,\%).
        \end{itemize}
  \item \textbf{Interface de démonstration (Streamlit) :}
        \begin{itemize}
          \item exploration du jeu de données : comptages, exemples par ville
                et par classe, transformations d'images ;
          \item \textbf{choix d'une image et d'un modèle, puis affichage du
                masque prédit} à côté du vrai masque, avec les scores sur
                l'image ;
          \item tableau de comparaison des trois modèles ;
          \item prise en compte des critères d'accessibilité WCAG essentiels.
        \end{itemize}
\end{enumerate}

\subsection*{Critère de succès}

EoMT dépasse la mIoU de la baseline de façon statistiquement significative,
avec un temps d'inférence compatible avec une utilisation via l'API du projet.

\end{document}
